\section{Equality Characterization}\label{sec:equality}
%==============================================================================


This section identifies the structural criterion under which the current upper theory is sharp. The upper theory is not presented only as a bound. The current equality criterion can be characterized exactly.

The original clique-fiber regime is now isolated as an equality subclass rather than only a threshold corollary. If a graph is generated by a surjective label map $x \mapsto b(x)$ with adjacency exactly between distinct vertices in the same fiber, then one representative from each fiber forms an independent set of size $|\mathcal B|$, while the complement graph is colorable by the same label map using $|\mathcal B|$ colors. This equality for cluster-graph structure is classical; the mechanized development packages the model-specific version needed here into a restricted-class theorem:
\[
C(G)=\vartheta_{\infty}(G)=\log |\mathcal B|.
\]
Thus the extended theory now contains both a genuine non-clique graph regime with strict upper/lower separation and a clean equality regime recovering the original fiber picture inside the same asymptotic Lov\'asz-$\vartheta$ framework. That equality is no longer confined to the abstract cluster-graph subclass. The current sharpness criterion is characterized exactly by transitivity of the base confusability relation:
\[
\text{confusability is transitive}
\iff
G_{\mathrm{conf}}=\mathrm{Cluster}(\Pi_{\mathrm{cc}}).
\]
Hence, if confusability is transitive, the base confusability graph is exactly the cluster graph on connected components, and the model-specific equality theorem yields
\[
C_\infty(\text{views})=\vartheta_\infty(\text{views})=\log |\Pi_{\mathrm{cc}}|,
\]
where $\Pi_{\mathrm{cc}}$ is the connected-component partition of the base confusability graph. Between arbitrary base models and full fiber coherence, a checkable sufficient condition is that the view family be \emph{meet-witnessed}: every pair of allowed views contains another allowed view inside their intersection. Under that condition, any two confusability witnesses can be pulled back to a common subview, so confusability is transitive and the same equality follows. Fiber coherence is then a stronger sufficient structural condition: if equality at one allowed location already fixes the full observation transcript, transitivity follows, the connected components are exactly the realized transcript fibers, and the equality sharpens further to
\[
C_\infty(\text{views})=\vartheta_\infty(\text{views})=\log |\mathcal T_{\mathrm{real}}|.
\]
At the witness level, this same collapse admits a local composition form: base confusability is the edge union of the single-view cluster graphs, and transitivity is equivalent to closure of those fixed-view witnesses under two-step composition through an intermediate state. This is the exact local criterion for the current equality mechanism: the collapse to a cluster graph is determined by how single-view witnesses compose, not by an opaque global coincidence of the full confusability graph.\leanmeta{\LHrng{MFT}{85}{102}, \LHrng{MFT}{107}{109}}

\medskip
\noindent\textbf{Hierarchy summary.} This equality route has one exact criterion and two stronger sufficient conditions. The table below records only the logical structure and the resulting capacity value.

\begingroup
\small
\setlength{\tabcolsep}{4pt}
\renewcommand{\arraystretch}{1.1}
\begin{center}
\begin{tabularx}{0.98\linewidth}{@{}>{\RaggedRight\arraybackslash}p{0.21\linewidth}>{\RaggedRight\arraybackslash}p{0.16\linewidth}>{\RaggedRight\arraybackslash}p{0.31\linewidth}>{\RaggedRight\arraybackslash}p{0.22\linewidth}@{}}
\toprule
\textbf{Condition} & \textbf{Necessity?} & \textbf{Structural consequence} & \textbf{Capacity value} \\
\midrule
Fiber coherence & No; stronger sufficient condition & Connected components are realized transcript fibers & $\log |\mathcal T_{\mathrm{real}}|$ \\
Meet-witness & No; sufficient for collapse & Base confusability collapses to the component cluster graph & $\log |\Pi_{\mathrm{cc}}|$ \\
Transitivity & Yes, for the current equality mechanism & $G_{\mathrm{conf}}=\mathrm{Cluster}(\Pi_{\mathrm{cc}})$ & $\log |\Pi_{\mathrm{cc}}|$ \\
\bottomrule
\end{tabularx}
\end{center}
\endgroup
